\subsubsection*{Statement L6c (Proved): first direct-image vanishing in every dimension}
\begin{theorem}
\label{B:L6-R1-vanishing}
Let $f:X\to Y$ be a proper birational morphism of integral
Noetherian schemes of characteristic $p>0$.  Suppose that $X$
is normal and $Y$ is regular.  Then
\[
                    R^1f_*\mathcal O_X=0.
\]
No dimension, projectivity, perfectness, or finite-type-over-a-field
assumption is required.  In particular, for proper varieties as in
B.1 or B.1g, pullback is an isomorphism on $H^1(\mathcal O)$ and
injective on $H^2(\mathcal O)$.  Thus the cohomological hypothesis
in L6a holds automatically for these morphisms.
\end{theorem}
\begin{proof}
We first record the semilinear fact needed over an imperfect residue
field.  Let $\kappa$ be separably closed of characteristic $p$, and
let $M\ne0$ be a finite-dimensional $\kappa$-vector space with an
injective additive map $F$ satisfying $F(av)=a^pF(v)$.  Put
\[
             V_n=\operatorname{span}_{\kappa}F^n(M).
\]
These are nonzero decreasing subspaces, and
$\operatorname{span}_{\kappa}F(V_n)=V_{n+1}$.  Choose $n$ after
the dimensions stabilize and put $V=V_n=V_{n+1}$.  In a basis
of $V$, the restriction of $F$ is $v\mapsto A v^{[p]}$ with
$A\in\operatorname{GL}_r(\kappa)$ and $r=\dim V>0$:
its column span is precisely $\operatorname{span}_{\kappa}F(V)=V$.
The kernel of $A v^{[p]}-v$ is finite \emph{\'etale} of degree
$p^r$.  Indeed, multiplication by $A^{-1}$ writes its equations
as $v_i^p-\sum_j(A^{-1})_{ij}v_j=0$; their pairwise relatively
prime leading monomials show that the monomials
$v_1^{e_1}\cdots v_r^{e_r}$ with $0\le e_i<p$ form a basis of
the quotient ring.  The Jacobian matrix of the original equations
is $-I_r$, so this finite scheme is \emph{\'etale}.  Since $\kappa$
is separably closed, it has $p^r>1$ rational points.  Consequently
$M$ has a nonzero element fixed by $F$.
This argument uses stabilized spans; injectivity of a semilinear
map alone does not imply that its matrix is invertible over an
imperfect field.

Suppose now that $\mathcal M=R^1f_*\mathcal O_X$ is nonzero.
It is coherent.  Choose a generic point of an irreducible
component of its support, localize $Y$ there, and then take a
strict henselization $A$.  The ring $A$ is Noetherian and regular,
and its residue field $\kappa$ is separably closed
\cite[Tags 06DI and 06LN]{StacksB}.  The resulting base change
$T\to\operatorname{Spec}A$ is proper and birational, and $T$
is normal: normality is preserved by localization and by the
ind-\'etale base change defining strict henselization.  It is
integral as well.  For example, flat base change preserves the
schematic density of the birational isomorphism locus, whose
base change is a nonempty open in the integral regular scheme
$\operatorname{Spec}A$; normality then excludes any additional
component.

By flat base change for coherent cohomology
\cite[Lemma 30.5.2, Tag 02KE]{StacksB},
\[
 H^0(T,\mathcal O_T)=A,\qquad
 M:=H^1(T,\mathcal O_T)\ne0,
\]
and $M$ is a finite-length $A$-module.  Finite length follows
from the choice of a minimal support prime; nonvanishing follows
from faithful flatness.  The identity in degree zero also follows
directly from proper birationality and normality of the base.

We apply B.1b to the finite flat groups $\alpha_p$ and
$(\mathbf Z/p\mathbf Z)_A$.  Its input applies exactly in this
scope: \cite[Proposition 3, p.~432, and Theorem 6,
pp.~433--434]{Marrama2026} require normality for full faithfulness
and regularity of the base, with complement of codimension at
least two, for torsor purity.  They impose no perfectness or
field-base hypothesis.  Hence
\[
 H^1_{\mathrm{fppf}}(A,G)
       \xrightarrow{\ \sim\ } H^1_{\mathrm{fppf}}(T,G_T)
       \quad(G=\alpha_p,\ \mathbf Z/p\mathbf Z).
\]
Quasi-coherent cohomology agrees in the Zariski and fppf
topologies \cite[Proposition 35.9.3, Tag 03DW]{StacksB}.
The sequence
\[
 0\longrightarrow\alpha_p\longrightarrow\mathbb G_a
       \xrightarrow{F}\mathbb G_a\longrightarrow0
\]
and the equality of global functions therefore identify
$H^1(A,\alpha_p)=A/F(A)$ with the same subgroup of
$H^1(T,\alpha_p)$.  The displayed torsor isomorphism forces
\[
                         \ker(F:M\to M)=0.
\]
Similarly, the Artin--Schreier sequence gives
\[
                      \ker(F-1:M\to M)=0,
\]
because every finite \emph{\'etale} torsor over the strictly
henselian local ring $A$ is trivial, so $H^1(A,\mathbf Z/p\mathbf Z)=0$.

Let $\mathfrak m$ be the maximal ideal of $A$ and choose $N$
with $\mathfrak m^N M=0$.  For $a\in\mathfrak m$, $v\in M$,
and $p^e\ge N$, semilinearity gives
\[
                    F^e(av)=a^{p^e}F^e(v)=0.
\]
The injectivity of $F$ implies $av=0$.  Thus $\mathfrak mM=0$,
so $M$ is a nonzero finite-dimensional $\kappa$-vector space.
The semilinear fact proved above now gives a nonzero fixed
point of $F$, contradicting $\ker(F-1)=0$.  This proves the
vanishing of $R^1f_*\mathcal O_X$.

Finally, the low-degree Leray sequence, together with
$f_*\mathcal O_X=\mathcal O_Y$, gives the asserted
$H^1$ isomorphism and $H^2$ injection.  L6a then supplies
the corresponding formal Picard isomorphism.
\end{proof}

\begin{remark}[Scope of the vanishing]
The proposed $R^1f_*\mathcal O_X\ne0$ counterexample route in
L6 is therefore impossible under its normal-source and
regular-target hypotheses in every dimension.  We have proved
only the first higher direct-image vanishing here; the theorem
does not assert that all $R^if_*\mathcal O_X$ vanish for a
normal source.  Nor does it settle the remaining imperfect-field
image problem in B.2, whose permitted models can have singular
targets.
\end{remark}
